%==============================================================================
\section{Data Preparation --- Chuẩn bị dữ liệu cho học máy}
%==============================================================================

\begin{dongco}[title={Bước then chốt ảnh hưởng mạnh tới mô hình}]
Chuẩn bị dữ liệu là bước quan trọng trong ML, ảnh hưởng lớn tới độ chính xác và hiệu quả mô hình cuối.
Cần chú ý chi tiết và hiểu sâu dữ liệu cùng bài toán.
\end{dongco}

\subsection{Data Preparation là gì?}

\begin{dinhnghia}[title={Data preparation}]
Quá trình xử lý dữ liệu thô: làm sạch, tổ chức và biến đổi để phù hợp thuật toán ML.
Đây là quá trình liên tục, ảnh hưởng lớn tới hiệu năng mô hình.
Dữ liệu sạch, có cấu trúc $\Rightarrow$ kết quả tốt hơn.
\end{dinhnghia}

\begin{vidu}[title={Sơ đồ quy trình chuẩn bị dữ liệu}]
\tsimg{10_du_lieu_trong_ml/machine_learning_data_preparation.jpg}
\end{vidu}

\subsection{Tầm quan trọng}

\begin{phuongphap}[title={Vì sao chuẩn bị dữ liệu quan trọng}]
\begin{itemize}
  \item \textbf{Cải thiện độ chính xác}: thuật toán học từ dữ liệu; dữ liệu sạch, có cấu trúc $\Rightarrow$ kết quả chính xác hơn.
  \item \textbf{Hỗ trợ feature engineering}: chọn/tạo feature mới --- chuẩn bị dữ liệu tạo nền cho bước này.
  \item \textbf{Chất lượng dữ liệu}: dữ liệu thu thập thường có lỗi, không nhất quán; làm sạch/biến đổi giúp rút insight và pattern.
  \item \textbf{Tăng tốc dự đoán}: dữ liệu đã chuẩn bị dễ phân tích và cho kết quả đáng tin hơn.
\end{itemize}
\end{phuongphap}

\subsection{Các bước trong quy trình chuẩn bị dữ liệu}

\begin{phuongphap}[title={Chuỗi bước target}]
Data Collection $\rightarrow$ Data Cleaning $\rightarrow$ Data Transformation $\rightarrow$ Data Reduction $\rightarrow$ Data Splitting.

target: dữ liệu chính xác, đầy đủ, liên quan cho phân tích và mô hình hoá.
\end{phuongphap}

\begin{luuy}[title={Không luôn tuần tự}]
Quy trình không nhất thiết theo thứ tự cố định (ví dụ có thể chia dữ liệu trước khi transform); có thể cần thu thập thêm dữ liệu.
\end{luuy}

\subsection{Data Collection và tích hợp}

\begin{ynghia}[title={Thu thập dữ liệu}]
Bước đầu: gom dữ liệu từ database, file văn bản, ảnh, âm thanh, web scraping, \ldots

Sau khi chọn dữ liệu thô, cần tiền xử lý để có insight --- đưa dữ liệu về dạng thuật toán ML dùng được.

Đôi khi thu thập đi kèm bước \textbf{data integration}: gộp dữ liệu từ nhiều nguồn (khớp/liên kết bản ghi, merge theo biến chung).
\end{ynghia}

\subsection{Data Cleaning}

\begin{dinhnghia}[title={Data cleaning}]
Nhận diện và sửa lỗi, giá trị thiếu, trùng lặp, outlier, \ldots
Đảm bảo dữ liệu chính xác, liên quan, không lỗi trước khi huấn luyện.
\end{dinhnghia}

\begin{phuongphap}[title={Các bước làm sạch}]
\begin{enumerate}
  \item \textbf{Xử lý trùng lặp}: nhận diện và xoá bằng \texttt{drop\_duplicates} trong Pandas.
  \item \textbf{Sửa lỗi cú pháp/cấu trúc}: chuẩn hoá định dạng, quy ước đặt tên.
  \item \textbf{Outlier}: phát hiện bằng z-score, IQR, hoặc ML (clustering, SVM, \ldots).
  \item \textbf{Missing values}: imputation (mean/median/mode, regression imputation, multiple imputation) hoặc deletion (mất dữ liệu).
  \item \textbf{Validation}: kiểm tra kiểu dữ liệu, mã, định dạng, khoảng giá trị trước khi lưu DB.
\end{enumerate}
\end{phuongphap}

\subsection{Data Transformation}

\begin{ynghia}[title={Biến đổi dữ liệu}]
Chuyển dữ liệu từ định dạng gốc sang định dạng phù hợp phân tích/mô hình: cấu trúc hoá, căn chỉnh, trích xuất, lưu trữ.
\end{ynghia}

\begin{phuongphap}[title={Kỹ thuật biến đổi phổ biến}]
Scaling, Normalization (L1, L2), Standardization, Binarization, Encoding, Log Transformation.
\end{phuongphap}

\subsubsection{Scaling (Min-Max)}

\begin{ynghia}[title={Rescaling}]
Thuộc tính thường có thang đo khác nhau; nhiều thuật toán ML cần rescale về cùng thang (thường 0--1).
Dùng \texttt{MinMaxScaler} của scikit-learn.
\end{ynghia}

\begin{vidu}[title={MinMaxScaler trên Pima Indians Diabetes}]
\begin{lstlisting}
from sklearn.datasets import fetch_openml
from numpy import set_printoptions
from sklearn import preprocessing

names = ['preg', 'plas', 'pres', 'skin', 'test', 'mass', 'pedi', 'age', 'class']
dataframe = fetch_openml("diabetes", version=1, as_frame=True, parser="auto").frame
dataframe.columns = names
array = dataframe.values

data_scaler = preprocessing.MinMaxScaler(feature_range=(0,1))
data_rescaled = data_scaler.fit_transform(array)

set_printoptions(precision=1)
print ("\nScaled data:\n", data_rescaled[0:10])
\end{lstlisting}
\textbf{Output} (10 hàng đầu): mọi giá trị nằm trong khoảng $[0,1]$.
\end{vidu}

\subsubsection{Normalization (L1 và L2)}

\begin{ynghia}[title={Normalization}]
Rescale sao cho mỗi \textbf{hàng} có độ dài 1 (hữu ích với sparse data nhiều số 0).
Dùng \texttt{Normalizer} của sklearn.
\end{ynghia}

\textbf{L1 Normalization}: tổng giá trị tuyệt đối mỗi hàng bằng 1 (Least Absolute Deviations).

\begin{vidu}[title={L1}]
\begin{lstlisting}
from sklearn.datasets import fetch_openml
from numpy import set_printoptions
from sklearn.preprocessing import Normalizer

names = ['preg', 'plas', 'pres', 'skin', 'test', 'mass', 'pedi', 'age', 'class']
dataframe = fetch_openml("diabetes", version=1, as_frame=True, parser="auto").frame
dataframe.columns = names
array = dataframe.values

Data_normalizer = Normalizer(norm='l1').fit(array)
Data_normalized = Data_normalizer.transform(array)

set_printoptions(precision=2)
print ("\nNormalized data:\n", Data_normalized [0:3])
\end{lstlisting}
\end{vidu}

\textbf{L2 Normalization}: tổng bình phương mỗi hàng bằng 1 (least squares).

\begin{vidu}[title={L2}]
\begin{lstlisting}
Data_normalizer = Normalizer(norm='l2').fit(array)
Data_normalized = Data_normalizer.transform(array)
set_printoptions(precision=2)
print ("\nNormalized data:\n", Data_normalized [0:3])
\end{lstlisting}
\end{vidu}

\subsubsection{Standardization}

\begin{ynghia}[title={Chuẩn hoá Gaussian}]
Biến đổi về phân phối Gaussian chuẩn: mean $= 0$, std $= 1$.
Hữu ích cho linear regression, logistic regression giả định phân phối Gaussian đầu vào.
Dùng \texttt{StandardScaler}.
\end{ynghia}

\begin{vidu}[title={StandardScaler}]
\begin{lstlisting}
from sklearn.datasets import fetch_openml
from sklearn.preprocessing import StandardScaler
from numpy import set_printoptions

names = ['preg', 'plas', 'pres', 'skin', 'test', 'mass', 'pedi', 'age', 'class']
dataframe = fetch_openml("diabetes", version=1, as_frame=True, parser="auto").frame
dataframe.columns = names
array = dataframe.values

data_scaler = StandardScaler().fit(array)
data_rescaled = data_scaler.transform(array)

set_printoptions(precision=2)
print ("\nRescaled data:\n", data_rescaled [0:5])
\end{lstlisting}
\end{vidu}

\subsubsection{Binarization}

\begin{ynghia}[title={Binarize / threshold}]
Chuyển dữ liệu thành nhị phân theo ngưỡng: trên ngưỡng $\rightarrow 1$, dưới $\rightarrow 0$.
Hữu ích khi có xác suất cần chuyển thành giá trị rời rạc.
Dùng \texttt{Binarizer}.
\end{ynghia}

\begin{vidu}[title={Binarizer với threshold = 0.5}]
\begin{lstlisting}
from sklearn.datasets import fetch_openml
from sklearn.preprocessing import Binarizer

names = ['preg', 'plas', 'pres', 'skin', 'test', 'mass', 'pedi', 'age', 'class']
dataframe = fetch_openml("diabetes", version=1, as_frame=True, parser="auto").frame
dataframe.columns = names
array = dataframe.values

binarizer = Binarizer(threshold=0.5).fit(array)
Data_binarized = binarizer.transform(array)

print ("\nBinary data:\n", Data_binarized [0:5])
\end{lstlisting}
\end{vidu}

\subsubsection{Encoding (Label Encoding)}

\begin{ynghia}[title={Label encoding}]
Sklearn thường yêu cầu nhãn dạng số thay vì chữ; chuyển nhãn chữ $\rightarrow$ số bằng \texttt{LabelEncoder}.
\end{ynghia}

\begin{vidu}[title={LabelEncoder}]
\begin{lstlisting}
import numpy as np
from sklearn import preprocessing

input_labels = ['red','black','red','green','black','yellow','white']

encoder = preprocessing.LabelEncoder()
encoder.fit(input_labels)

test_labels = ['green','red','black']
encoded_values = encoder.transform(test_labels)
print("\nLabels =", test_labels)
print("Encoded values =", list(encoded_values))

encoded_values = [3,0,4,1]
decoded_list = encoder.inverse_transform(encoded_values)
print("\nEncoded values =", encoded_values)
print("\nDecoded labels =", list(decoded_list))
\end{lstlisting}
\textbf{Output}:
\begin{lstlisting}
Labels = ['green', 'red', 'black']
Encoded values = [1, 2, 0]
Encoded values = [3, 0, 4, 1]
Decoded labels = ['white', 'black', 'yellow', 'green']
\end{lstlisting}
One-hot và target encoding chi tiết hơn ở bài \textbf{Categorical Data}.
\end{vidu}

\subsubsection{Log Transformation}

\begin{ynghia}[title={Log transform}]
Thường dùng với dữ liệu lệch (skewed): áp logarit tự nhiên lên giá trị để thay đổi thang đo.
\end{ynghia}

\subsection{Data Reduction}

\begin{dinhnghia}[title={Giảm dữ liệu}]
Chọn tập con feature hoặc quan sát liên quan nhất để giảm kích thước dataset, giảm nhiễu, có thể cải thiện mô hình.
Hữu ích khi dataset rất lớn hoặc nhiều phần không liên quan.
\end{dinhnghia}

\begin{phuongphap}[title={Kỹ thuật}]
\begin{itemize}
\item \textbf{Dimensionality reduction} (giảm chiều không mất thông tin quan trọng)
\item \textbf{Discretization} (rời rạc hoá biến liên tục như thời gian, nhiệt độ).
\end{itemize}
\end{phuongphap}

\subsection{Data Splitting}

\begin{dinhnghia}[title={Chia dữ liệu}]
Bước cuối chuẩn bị dữ liệu cho ML:
\begin{itemize}
  \item \textbf{Training}: tập huấn luyện, học pattern.
  \item \textbf{Validation}: đánh giá trong lúc huấn luyện.
  \item \textbf{Testing}: đánh giá mô hình đã huấn luyện.
\end{itemize}
\end{dinhnghia}

\subsection{Ví dụ Python: breast cancer dataset}

\begin{vidu}[title={load, split, StandardScaler}]
\begin{lstlisting}
from sklearn.datasets import load_breast_cancer
from sklearn.model_selection import train_test_split
from sklearn.preprocessing import StandardScaler

# load the dataset
data = load_breast_cancer()

# separate the features and target
X = data.data
y = data.target

# split the data into training and testing sets
X_train, X_test, y_train, y_test = train_test_split(X, y, test_size=0.2, random_state=42)

# normalize the data using StandardScaler
scaler = StandardScaler()
X_train = scaler.fit_transform(X_train)
X_test = scaler.transform(X_test)
\end{lstlisting}
Nạp dataset, tách feature/target, chia train/test (80/20), chuẩn hoá bằng \texttt{StandardScaler} (\texttt{fit} trên train, \texttt{transform} trên test).
Quan trọng cho SVM, neural networks khi các feature cần cùng thang đo.
\end{vidu}

\subsection{Data Preparation và Feature Engineering}

\begin{ynghia}[title={Hai mặt của tiền xử lý}]
Feature engineering tạo feature mới từ dữ liệu hiện có (kết hợp, biến đổi, kiến thức miền).
Data preparation và feature engineering đi song song trong pipeline tiền xử lý tổng thể.
\end{ynghia}
